\documentclass[11pt]{article}

\usepackage[a4paper,margin=25mm]{geometry}
\usepackage{amsmath,amssymb,mathtools}
\usepackage{booktabs,longtable,array}
\usepackage[hidelinks]{hyperref}
\usepackage[T1]{fontenc}

\newcommand{\LSE}{\operatorname{LSE}_2}
\newcommand{\Anc}{\operatorname{Anc}}
\newcommand{\Sub}{\operatorname{Sub}}
\newcommand{\dd}{\mathrm{d}}
\newcommand{\kernel}[1]{\nolinkurl{#1}}
\newcommand{\kernelref}[2]{\href{../../gpurec/core/kernels/#1}{\nolinkurl{#2}}}
\newcommand{\R}{\mathbb{R}}

\title{Mathematical Reference for the \texttt{gpurec} Triton Kernels}
\author{gpurec developers}
\date{\today}

\begin{document}
\maketitle
\tableofcontents

\section{Scope}

This document is the canonical mathematical description of the device kernels
in \path{gpurec/core/kernels}.  Python docstrings state interfaces and point
here; equations, invariants, and relationships between primal, tangent,
adjoint, and second-order kernels live here.

The implementation uses base-2 logarithms.  Define
\begin{equation}
  \LSE(x_1,\ldots,x_n)
  = \log_2\!\left(\sum_{k=1}^n 2^{x_k}\right).
  \label{eq:lse}
\end{equation}
Impossible events have log-weight $-\infty$.

\section{Notation and model objects}

Let $\mathcal S$ be the species-tree states, $S=|\mathcal S|$, and let
$c_1(s)$, $c_2(s)$, and $\operatorname{pa}(s)$ denote the two children and
parent of state $s$.  We include $s$ itself in its ancestral path:
\[
  \Anc(s)=\{s,\operatorname{pa}(s),\operatorname{pa}^2(s),\ldots\}.
\]
Likewise, $\Sub(s)$ is the species subtree rooted at $s$, including $s$.
The valid recipients for a transfer whose donor is $s$ are
\begin{equation}
  \mathcal V(s)=\mathcal S\setminus\Anc(s).
  \label{eq:valid-recipients}
\end{equation}
Thus neither the donor itself nor one of its ancestors can receive the
transfer.

In the weighted-receiver model, let $\omega_r\in\R$ be the unconstrained
receiver logit.  Its normalized base-2 log-probability is
\begin{equation}
  \rho_r=\log_2\!\left(
    \frac{\exp(\omega_r)}{\sum_{t\in\mathcal S}\exp(\omega_t)}
  \right).
  \label{eq:receiver-log-softmax}
\end{equation}
Adding a common constant to every $\omega_r$ therefore leaves the model
unchanged.  The kernels receive $\rho$ after this normalization and use the
effective receiver log-measure
\begin{equation}
  \alpha_r=
  \begin{cases}
    \rho_r,&\text{weighted-receiver branch},\\
    0,&\text{uniform-receiver branch}.
  \end{cases}
  \label{eq:receiver-measure}
\end{equation}
The uniform branch therefore sums one unit of mass per valid recipient.  It
does not load or differentiate a receiver-weight vector, which is why retaining
the compile-time branch avoids work.

For gene family $g$ and species state $s$, write
$\ell^S_{g,s}$, $\ell^D_{g,s}$, and $\ell^L_{g,s}$ for log-probabilities of
speciation, duplication, and loss.  A parameter can be shared or vary along
the family and species axes; this changes storage only, not the equations.
Write $\ell^T_{g,s}$ for the transfer log-probability.  The donor-specific
transfer base term is
\begin{equation}
  m_g(s)=\ell^T_g(s)
  -\log_2\!\left(\sum_{r\in\mathcal V(s)}2^{\alpha_r}\right).
  \label{eq:max-transfer}
\end{equation}
It is named \texttt{max\_transfer} in kernel interfaces.  In particular,
$m_g(s)=\ell^T_g(s)-\log_2|\mathcal V(s)|$ in the uniform branch.  Parameter
extraction constructs $m$ and $\rho$ together according to
Equation~\eqref{eq:max-transfer}; kernel boundaries accept them independently,
so their JVPs and VJPs are also exposed independently before the outer chain
rule couples them again.

$E_g(s)$ is the log extinction probability.  For gene-clade row $i$,
$\Pi_i(s)$ is its reconciliation log-likelihood and $\bar\Pi_i(s)$ is the
transfer-complement likelihood.  A row may be stored as a residual plus a
scalar gauge:
\begin{equation}
  \Pi_i^{\mathrm{abs}}(s)=\Pi_i^{\mathrm{res}}(s)+o_i,
  \qquad
  \bar\Pi_i^{\mathrm{abs}}(s)=\bar\Pi_i^{\mathrm{res}}(s)+\bar o_i.
  \label{eq:gauge}
\end{equation}
All probabilities and derivatives below refer to the represented absolute
quantity.  A tangent of a stored residual is the tangent of that absolute
quantity with the chosen gauge held fixed; selecting a different gauge does
not change the mathematical derivative.

\section{Reusable log-sum-exp and transfer complement}

For any log-vector $x(s)$, the transfer-complement operator is
\begin{equation}
  \mathcal B[x](s)
  =m(s)+\log_2\!\left(
    \sum_{r\in\mathcal V(s)}2^{\alpha_r+x(r)}
  \right).
  \label{eq:transfer-complement-direct}
\end{equation}
Equation~\eqref{eq:max-transfer} makes
$\mathcal B[0](s)=\ell^T(s)$: the receiver measure redistributes a fixed
transfer probability among the valid recipients.

For stable evaluation, define
\begin{align}
  M_x &= \max_{r\in\mathcal S}\{\alpha_r+x(r)\}, \\
  q_x(r)&=2^{\alpha_r+x(r)-M_x},\\
  R_x &= \sum_{r\in\mathcal S}q_x(r), \\
  A_x(s) &= \sum_{r\in\Anc(s)}q_x(r),\\
  D_x(s)&=R_x-A_x(s)
         =\sum_{r\in\mathcal V(s)}q_x(r).
\end{align}
The mass outside the donor's ancestral lineage is
\begin{equation}
  \mathcal B[x](s)
  = m(s)+M_x+\log_2 D_x(s).
  \label{eq:transfer-complement}
\end{equation}
It is $-\infty$ when the complementary mass is zero.  Kernel helpers
\kernelref{pi_forward.py}{_compute_total_receiver_mass} and
\kernelref{pi_forward.py}{_compute_transfer_complement} evaluate these quantities.
The same operator produces $\bar E$ from $E$ and $\bar\Pi$ from $\Pi$.

The row maximum $M_x$ is only a numerical gauge.  Equation
\eqref{eq:transfer-complement-direct} contains no $M_x$, so it is exact to hold
$M_x$ fixed in every JVP and VJP.  Equivalently, differentiating $M_x$ would
add $\dd M_x$ explicitly and subtract the same quantity through
$\dd\log_2 D_x(s)$.

\section{Extinction fixed point}

For each family and species state, define
\begin{align}
  e^{\mathrm{spec}}(s) &= \ell^S(s)+E(c_1(s))+E(c_2(s)), \\
  e^{\mathrm{dup}}(s) &= \ell^D(s)+2E(s), \\
  e^{\mathrm{trans}}(s) &= E(s)+\bar E(s), \\
  e^{\mathrm{loss}}(s) &= \ell^L(s), \\
  \bar E(s) &= \mathcal B[E](s).
\end{align}
One E-step is
\begin{equation}
  F_E(E)(s)=\LSE\!\left(
    e^{\mathrm{spec}}(s),e^{\mathrm{dup}}(s),
    e^{\mathrm{trans}}(s),e^{\mathrm{loss}}(s)
  \right).
  \label{eq:e-step}
\end{equation}
The fixed point satisfies $E^*=F_E(E^*)$.

\subsection{E-step tangent}

For $z=\LSE(a_1,\ldots,a_n)$, define event probabilities
\begin{equation}
  w_k=2^{a_k-z}, \qquad \sum_k w_k=1.
  \label{eq:lse-weights}
\end{equation}
Its directional derivative is
\begin{equation}
  \dd z=\sum_k w_k\,\dd a_k.
  \label{eq:lse-jvp}
\end{equation}
For the transfer complement, let $q_E(r)$ and $D_E(s)$ be the masses defined
above.  With the stabilizing maximum frozen,
\begin{equation}
  \dd\bar E(s)
  =\dd m(s)
   +\frac{
     \sum_{r\in\mathcal V(s)}q_E(r)
       (\dd\alpha_r+\dd E(r))
   }{D_E(s)}.
  \label{eq:ebar-jvp}
\end{equation}
Here $\dd\alpha=\dd\rho$ in the weighted branch and $\dd\alpha=0$ in the
uniform branch.
The tangent fixed point solves
\begin{equation}
  (I-J_E)\,\dd E^*=J_p\,\dd p,
  \label{eq:e-tangent-fixed-point}
\end{equation}
where $J_E=\partial F_E/\partial E$ and $J_p$ contains parameter
derivatives.  The extinction JVP kernel evaluates one application of the
right-hand side of this linearized recurrence; Section~\ref{sec:kernel-index}
links it to its source file.

\subsection{E-step adjoint and curvature}

For an output cotangent $g(s)$, the log-sum-exp event cotangent is
\begin{equation}
  q_k(s)=g(s)2^{e_k(s)-F_E(E)(s)}.
  \label{eq:e-adjoint-event}
\end{equation}
Let $g_{\bar E}(s)$ be any additional cotangent arriving at the materialized
$\bar E$ output.  The total cotangent of $\mathcal B[E](s)$ is
$b(s)=q_{\mathrm{trans}}(s)+g_{\bar E}(s)$.  Define
\begin{equation}
  u(s)=\frac{b(s)}{D_E(s)}.
\end{equation}
Then $m(s)$ receives $b(s)$, while each receiver state $r$ contributes the
following amount to the cotangent of $E(r)$:
\begin{equation}
  q_E(r)\left(
    \sum_{s\in\mathcal S}u(s)-\sum_{s\in\Sub(r)}u(s)
  \right).
  \label{eq:complement-vjp}
\end{equation}
The same quantity is the cotangent of $\rho_r$ in the weighted branch.  There
is no receiver-weight cotangent in the uniform branch.  The identity follows
because $r\notin\mathcal V(s)$ exactly when $s\in\Sub(r)$.
The prepare and finalize backward kernels split this formula before and after
the ancestral exclusion accumulation.

For an objective cotangent $b_E$ arriving at the converged extinction state,
the fixed-point adjoint is the solution of
\begin{equation}
  (I-J_E^T)w_E=b_E.
  \label{eq:e-adjoint-fixed-point}
\end{equation}
One E-step VJP supplies applications of $J_E^T$ inside this solve.  After the
solve, applying the same VJP with $w_E$ yields $J_p^T w_E$, which is combined
with any explicit parameter cotangents.  Receiver-log-probability and
\texttt{max\_transfer} cotangents remain separate at the kernel boundary; the
outer differentiation through Equations~\eqref{eq:receiver-log-softmax} and
\eqref{eq:max-transfer} maps both contributions back to the raw receiver
logits.

At fixed $g$, the second-order contraction needed by an exact
Hessian-vector product is
\begin{equation}
  D_x[J_{F_E}(x)^Tg][\dd x].
  \label{eq:e-so}
\end{equation}
More generally, if $z=\LSE(a_1,\ldots,a_n)$, $w_k=2^{a_k-z}$, and
$q_k=g w_k$ is an event VJP with $g$ held fixed, then
\begin{equation}
  \dd q_k=\ln(2)q_k(\dd a_k-\dd z).
  \label{eq:event-weight-derivative}
\end{equation}
For the complement part, put
\begin{align}
  D_E(s)&=\sum_{r\in\mathcal V(s)}q_E(r), &
  b(s)&=q_{\mathrm{trans}}(s)+g_{\bar E}(s), &
  u(s)&=\frac{b(s)}{D_E(s)}.
\end{align}
At fixed $g_{\bar E}$,
\begin{align}
  \dd D_E(s)
    &=\ln(2)\sum_{r\in\mathcal V(s)}q_E(r)
      (\dd\alpha_r+\dd E(r)), \\
  \dd b(s)&=\dd q_{\mathrm{trans}}(s).
\end{align}
The quotient rule gives
\begin{equation}
  \dd u(s)=\frac{\dd b(s)-u(s)\,\dd D_E(s)}{D_E(s)}.
  \label{eq:quotient-tangent}
\end{equation}
Equations~\eqref{eq:event-weight-derivative} and
\eqref{eq:quotient-tangent}, followed by the product rule applied to
Equation~\eqref{eq:complement-vjp}, define the E-step second-order kernels.

\section{Within-wave reconciliation recurrence}

For gene-clade row $i$ and species state $s$, the local event terms are
formed from
\begin{align}
  \ell^{DL}(s)&=1+\ell^D(s)+E(s), \\
  \ell^{SL_1}(s)&=\ell^S(s)+E(c_2(s)), \\
  \ell^{SL_2}(s)&=\ell^S(s)+E(c_1(s)).
\end{align}
The $1=\log_2 2$ in $\ell^{DL}$ accounts for the two choices of which
duplicate lineage is lost.  The extinction factor paired with one reconciled
species child is the extinction probability of the other species child.  Let
$\lambda(i)$ be the observed species of a leaf gene clade.  The six named
local terms are
\begin{align}
  p^{\mathrm{dup\text{-}loss}}(s)
    &=\ell^{DL}(s)+\Pi_i(s), \\
  p^{\mathrm{trans\text{-}loss}}(s)
    &=\Pi_i(s)+\bar E(s), \\
  p^{\mathrm{trans}}(s)
    &=\bar\Pi_i(s)+E(s), \\
  p^{\mathrm{spec},c_1}(s)
    &=\ell^{SL_1}(s)+\Pi_i(c_1(s)), \\
  p^{\mathrm{spec},c_2}(s)
    &=\ell^{SL_2}(s)+\Pi_i(c_2(s)), \\
  p^{\mathrm{leaf}}(s)
    &=
    \begin{cases}
      \ell^S(s),&s=\lambda(i),\\
      -\infty,&s\ne\lambda(i).
    \end{cases}
\end{align}
The \texttt{leaf\_logp} passed by the reconciliation solver is therefore the
same speciation log-probability $\ell^S$; the separate name describes its role
as the leaf source.

Collect the six within-wave terms in the local event aggregate
\begin{equation}
  L_i(s)=\LSE\bigl(
    p^{\mathrm{dup\text{-}loss}},
    p^{\mathrm{trans\text{-}loss}},
    p^{\mathrm{trans}},
    p^{\mathrm{spec},c_1},
    p^{\mathrm{spec},c_2},
    p^{\mathrm{leaf}}
  \bigr)(s).
  \label{eq:local-event-aggregate}
\end{equation}
If a gene clade has the split aggregate $D_i(s)$, its update is
\begin{equation}
  F_\Pi(\Pi_i)(s)
  =\LSE\bigl(L_i(s),D_i(s)\bigr),
  \label{eq:pi-wave}
\end{equation}
with an absent split or leaf term set to $-\infty$.  The transfer complement is
\begin{equation}
  \bar\Pi_i(s)=\mathcal B[\Pi_i](s).
  \label{eq:pibar}
\end{equation}

It is useful to factor a local event's probability into
\begin{align}
  \beta_{i,e}(s)&=2^{p^e_i(s)-L_i(s)},
    &\sum_e\beta_{i,e}(s)&=1,\\
  \gamma_i(s)&=2^{L_i(s)-F_\Pi(\Pi_i)(s)}.
  \label{eq:within-wave-probability}
\end{align}
Thus the probability of local event $e$ in the complete recurrence is
\begin{equation}
  a_{i,e}(s)=\gamma_i(s)\beta_{i,e}(s)
  =2^{p^e_i(s)-F_\Pi(\Pi_i)(s)}.
  \label{eq:complete-local-event-probability}
\end{equation}
Backward and second-order kernels use this factorization; the tangent kernel
may form the same complete probability directly.

\subsection{Equation-aligned implementation vocabulary}

The implementation spells out the event role rather than using numeric
buckets.  For each event name \texttt{event}, the identifiers have a uniform
meaning:
\begin{center}
\begin{tabular}{ll}
\toprule
Identifier suffix & Mathematical quantity \\
\midrule
\texttt{\_log\_term} & the corresponding $p^{\mathrm{event}}(s)$ \\
\texttt{\_mass} & $2^{p^{\mathrm{event}}(s)-a(s)}$ for the local numerical maximum $a(s)$ \\
\texttt{\_probability} & either $\beta_{i,e}$ or the explicitly named complete recurrence probability \\
\texttt{\_event\_vjp} & $v(s)\gamma_i(s)\beta_{i,e}(s)$ \\
\texttt{d\_} prefix & directional derivative of the named quantity \\
\bottomrule
\end{tabular}
\end{center}
Here \texttt{event} is one of the following equation-aligned prefixes:
\begin{center}
\begin{tabular}{lll}
\texttt{duplication\_loss} & \texttt{transfer\_loss} & \texttt{transfer} \\
\texttt{speciation\_child1} & \texttt{speciation\_child2} & \texttt{leaf\_observation}
\end{tabular}
\end{center}
The combined
\texttt{speciation\_leaf\_event\_vjp} is the sum of the last three VJPs
because all three depend on the speciation parameter.  This vocabulary is
used unchanged by the primal, JVP, VJP, and VJP directional-derivative
kernels.

\kernelref{pi_forward.py}{_initialize_leaf_reconciliation_likelihood_kernel} evaluates
Equation~\eqref{eq:pi-wave} for leaf rows from their observation source.
\kernelref{pi_forward.py}{_update_reconciliation_likelihood_kernel}
evaluates the general recurrence and publishes the final row gauges.

The wave tangent applies Equations~\eqref{eq:lse-jvp} and
\eqref{eq:ebar-jvp} to Equations~\eqref{eq:local-event-aggregate}--
\eqref{eq:pibar}.  The self-loop variant iterates or exposes the linear map
induced by the dependence of the new row on the old row.

For one clade row, let
$J_{i,\mathrm{self}}=\partial F_\Pi(\Pi_i)/\partial\Pi_i$ while holding child
clade rows, $D_i$, extinction quantities, and parameters fixed.  Define the
receiver probability conditional on donor $s$ by
\begin{equation}
  r_s(t)=
  \begin{cases}
    q_{\Pi_i}(t)/D_{\Pi_i}(s),&t\in\mathcal V(s),\\
    0,&t\notin\mathcal V(s).
  \end{cases}
  \label{eq:conditional-receiver-probability}
\end{equation}
For a row tangent $h$, the within-row Jacobian is therefore
\begin{align}
  (J_{i,\mathrm{self}}h)(s)
  ={}&\bigl(a_{i,\mathrm{dup\text{-}loss}}(s)
           +a_{i,\mathrm{trans\text{-}loss}}(s)\bigr)h(s) \notag\\
    &+a_{i,\mathrm{trans}}(s)
      \sum_t r_s(t)h(t) \notag\\
    &+a_{i,\mathrm{spec},c_1}(s)h(c_1(s))
     +a_{i,\mathrm{spec},c_2}(s)h(c_2(s)).
  \label{eq:wave-self-jacobian}
\end{align}
The leaf and gene-split terms have no within-row derivative.  Parameter and
receiver-weight tangents enter the full wave JVP as additional source terms;
they are not part of $J_{i,\mathrm{self}}h$.

The wave adjoint solves
\begin{equation}
  (I-J_{i,\mathrm{self}}^T)v_i=b_i,
  \label{eq:wave-adjoint-solve}
\end{equation}
where one application of the transpose is
\begin{align}
  (J_{i,\mathrm{self}}^Tv)(t)
  ={}&v(t)\bigl(a_{i,\mathrm{dup\text{-}loss}}(t)
                    +a_{i,\mathrm{trans\text{-}loss}}(t)\bigr) \notag\\
    &+\sum_{s:c_1(s)=t}v(s)a_{i,\mathrm{spec},c_1}(s)
     +\sum_{s:c_2(s)=t}v(s)a_{i,\mathrm{spec},c_2}(s) \notag\\
    &+q_{\Pi_i}(t)\left(
       \sum_s u_i(s)-\sum_{s\in\Sub(t)}u_i(s)
      \right),
  \label{eq:wave-self-transpose}
\end{align}
with
\begin{equation}
  u_i(s)=\frac{v(s)a_{i,\mathrm{trans}}(s)}{D_{\Pi_i}(s)}.
\end{equation}
This is exactly the diagonal, child-edge, and transfer-subtree decomposition
materialized by the self-loop VJP preparation kernel and consumed by the
transpose kernel.

After solving, the VJP of a local event log-term is simply
\begin{equation}
  h_{i,e}(s)=v(s)a_{i,e}(s).
  \label{eq:wave-event-vjp}
\end{equation}
Consequently, the rate-term VJPs are
\begin{align}
  \ell^D &: h_{\mathrm{dup\text{-}loss}},\\
  m &: h_{\mathrm{trans}},\\
  \ell^S &: h_{\mathrm{spec},c_1}
    +h_{\mathrm{spec},c_2}+h_{\mathrm{leaf}}.
\end{align}
The extinction quantities receive
\begin{align}
  E(s)&:\ h_{\mathrm{dup\text{-}loss}}(s)+h_{\mathrm{trans}}(s),\\
  \bar E(s)&:\ h_{\mathrm{trans\text{-}loss}}(s),\\
  E(c_2(s))&:\ h_{\mathrm{spec},c_1}(s),\\
  E(c_1(s))&:\ h_{\mathrm{spec},c_2}(s).
\end{align}
The $m$ VJP is named \texttt{max\_transfer} in the interfaces.  In the
weighted branch, the same transfer-subtree term in
Equation~\eqref{eq:wave-self-transpose} is also the VJP of $\rho_t$.
Cross-clade derivatives through $D_i$ are not part of this self-loop solve;
the gene-split VJP kernels scatter them afterward.

At fixed $v$, differentiating Equation~\eqref{eq:wave-event-vjp} gives
\begin{equation}
  \dd h_{i,e}(s)
  =v(s)\ln(2)a_{i,e}(s)
    \bigl(\dd p_i^e(s)-\dd F_\Pi(\Pi_i)(s)\bigr).
  \label{eq:wave-event-so}
\end{equation}
The wave second-order kernel computes all such event-VJP derivatives together
with
\begin{equation}
  D_x[J_{i,\mathrm{self}}(x)^Tv_i][\dd x],
  \label{eq:wave-so}
\end{equation}
using the product rule for the transfer tree.  Here $x$ includes the primal
row, extinction quantities, rate terms, split aggregate, and---in the weighted
branch---receiver log-probabilities.  The supplied $v_i$ is held fixed; the
separate tangent-adjoint solve accounts for $\dd v_i$.

\section{Duplication-transfer-speciation split reduction}

Consider a gene split $n$ with left child row $L_n$, right child row $R_n$,
parent row $P_n$, and split log-prior $\sigma_n$.  For species state $s$, let
\begin{align}
  d_{\mathrm{dup}}^{(n)}(s)&=\sigma_n+\ell^D(s)+\Pi_{L_n}(s)+\Pi_{R_n}(s), \\
  d_{\mathrm{trans},L}^{(n)}(s)&=\sigma_n+\Pi_{L_n}(s)+\bar\Pi_{R_n}(s), \\
  d_{\mathrm{trans},R}^{(n)}(s)&=\sigma_n+\Pi_{R_n}(s)+\bar\Pi_{L_n}(s), \\
  d_{\mathrm{spec},LR}^{(n)}(s)&=\sigma_n+\ell^S(s)
    +\Pi_{L_n}(c_1(s))+\Pi_{R_n}(c_2(s)), \\
  d_{\mathrm{spec},RL}^{(n)}(s)&=\sigma_n+\ell^S(s)
    +\Pi_{R_n}(c_1(s))+\Pi_{L_n}(c_2(s)).
\end{align}
These are respectively duplication, transfer with the left child retained,
transfer with the right child retained, and the two speciation assignments.
The parent split contribution is
\begin{equation}
  D_P(s)=\underset{n:\,P_n=P}{\operatorname{LSE}_2}
  \bigl(
    d_{\mathrm{dup}}^{(n)},d_{\mathrm{trans},L}^{(n)},
    d_{\mathrm{trans},R}^{(n)},d_{\mathrm{spec},LR}^{(n)},
    d_{\mathrm{spec},RL}^{(n)}
  \bigr)(s).
  \label{eq:dts}
\end{equation}
The corresponding code prefixes are
\begin{center}
\begin{tabular}{lll}
\texttt{duplication} & \texttt{transfer\_left\_retained} & \texttt{transfer\_right\_retained} \\
\texttt{speciation\_lr} & \texttt{speciation\_rl} &
\end{tabular}
\end{center}
The suffixes \texttt{log\_term},
\texttt{probability}, and \texttt{event\_vjp} have the same meanings as in
the within-wave vocabulary above.
Gauge-frame shifts from Equation~\eqref{eq:gauge} are added before forming
event probabilities in residual storage.  Equations~\eqref{eq:dts} and all
probabilities below are equations for the represented absolute values; the
frame shifts only recover those values without materializing them.

The single-split kernel evaluates Equation~\eqref{eq:dts} directly.  For a
parent with multiple splits, a staging kernel forms partial stabilized sums
and a finalization kernel merges them.  All three are linked in
Section~\ref{sec:kernel-index}.

\subsection{DTS tangent and adjoint}

With
\begin{equation}
  w_k^{(n)}(s)=2^{d_k^{(n)}(s)-D_P(s)},
\end{equation}
the five event tangents are
\begin{align}
  \dd d_{\mathrm{dup}}^{(n)}
    &=\dd\ell^D+\dd\Pi_{L_n}+\dd\Pi_{R_n},\\
  \dd d_{\mathrm{trans},L}^{(n)}
    &=\dd\Pi_{L_n}+\dd\bar\Pi_{R_n},\\
  \dd d_{\mathrm{trans},R}^{(n)}
    &=\dd\Pi_{R_n}+\dd\bar\Pi_{L_n},\\
  \dd d_{\mathrm{spec},LR}^{(n)}
    &=\dd\ell^S+\dd\Pi_{L_n}(c_1)+\dd\Pi_{R_n}(c_2),\\
  \dd d_{\mathrm{spec},RL}^{(n)}
    &=\dd\ell^S+\dd\Pi_{R_n}(c_1)+\dd\Pi_{L_n}(c_2).
  \label{eq:dts-event-jvps}
\end{align}
The split prior is static metadata and therefore has no tangent.  Then
the DTS tangent is
\begin{equation}
  \dd D_P(s)=
  \sum_{n:\,P_n=P}\sum_{k=0}^4w_k^{(n)}(s)\,\dd d_k^{(n)}(s).
  \label{eq:dts-jvp}
\end{equation}
Equation~\eqref{eq:dts-jvp} describes the standalone split-reduction JVP.  Its
normalizer is $D_P$ because its output is exactly $D_P$.

The split VJP used by the complete reconciliation recurrence has a different
normalizer.  A raw split event $d_k^{(n)}$ is one term of
$\Pi_P=F_\Pi(\Pi_P)$, so define
\begin{equation}
  \eta_k^{(n)}(s)=2^{d_k^{(n)}(s)-\Pi_P(s)}.
  \label{eq:dts-full-probability}
\end{equation}
For a solved parent cotangent $v_P(s)$, the event VJP is
$h_k^{(n)}(s)=v_P(s)\eta_k^{(n)}(s)$.  The gene-split VJP kernel scatters
these quantities to $L_n$, $R_n$, and the event parameters: duplication
contributes to both child rows and $\ell^D$; each transfer event contributes
to its retained $\Pi$ row and the other child's $\bar\Pi$ row; and the two
speciation assignments scatter to their selected species children and
$\ell^S$.  Both transfer events also contribute to $m$ through their
$\bar\Pi$ operand.  In general the $\eta$ values sum to the probability
assigned to the split aggregate inside
Equation~\eqref{eq:pi-wave}, not to one.  This distinction between $D_P$ in
Equation~\eqref{eq:dts-jvp} and $\Pi_P$ in
Equation~\eqref{eq:dts-full-probability} is essential.

Transfer terms contain $\bar\Pi$, whose adjoint has a convenient tree form.
For one child-side complement $\bar\Pi_X$, let $b_X(s)$ be the transfer-event
VJP arriving at donor $s$, and define
\begin{equation}
  u_X(s)=\frac{b_X(s)}{D_{\Pi_X}(s)},\qquad
  A_X=\sum_j u_X(j), \qquad
  U_{X,r}=\sum_{s\in\Sub(r)}u_X(s), \qquad
  p_{X,r}=2^{\alpha_r+\Pi_X(r)-M_{\Pi_X}}.
\end{equation}
The contribution to receiver state $r$ is
\begin{equation}
  C_{X,r}=p_{X,r}(A_X-U_{X,r}).
  \label{eq:tree-complement}
\end{equation}
This is Equation~\eqref{eq:complement-vjp} written as a bottom-up subtree
reduction.  In the weighted branch $C_{X,r}$ also contributes to the cotangent of
$\rho_r$; in the uniform branch no such parameter exists.
Compact level tables contain internal species-tree nodes in bottom-up order,
so every $U_{X,r}$ is formed after its child sums.

\subsection{DTS second-order contraction}

At fixed parent cotangent $v_P$, differentiating the complete-recurrence split
event VJP gives
\begin{equation}
  \dd h_k^{(n)}=v_P\ln(2)\eta_k^{(n)}
  (\dd d_k-\dd\Pi_P).
  \label{eq:dts-event-so}
\end{equation}
The subtraction is $\dd\Pi_P$, not $\dd D_P$, for the same reason as in
Equation~\eqref{eq:dts-full-probability}.
For a transfer side, the donor coefficient also changes:
\begin{align}
  \dd D_{\Pi_X}(s)
    &=\ln(2)\sum_{r\in\mathcal V(s)}p_{X,r}
      (\dd\alpha_r+\dd\Pi_X(r)),\\
  \dd u_X(s)
    &=\frac{\dd b_X(s)-u_X(s)\dd D_{\Pi_X}(s)}{D_{\Pi_X}(s)}.
  \label{eq:dts-donor-so}
\end{align}
The split-event second-order kernel stages $u_X$ and $\dd u_X$.
Differentiating Equation~\eqref{eq:tree-complement} then gives
\begin{equation}
  \dd C_{X,r}=(\dd p_{X,r})(A_X-U_{X,r})
    +p_{X,r}(\dd A_X-\dd U_{X,r}),
  \label{eq:tree-complement-so}
\end{equation}
where
\begin{equation}
  \dd p_{X,r}=\ln(2)p_{X,r}
    (\dd\alpha_r+\dd\Pi_X(r))
\end{equation}
with the stabilizing row maximum frozen.  The split second-order kernel
implements Equation~\eqref{eq:dts-event-so}; the tree second-order kernel
implements Equation~\eqref{eq:tree-complement-so}.

\section{Kernel index}
\label{sec:kernel-index}

The following table accounts for every \texttt{@triton.jit} device function in
\texttt{gpurec/core/kernels}.  The first three entries are inline helpers used
inside launch kernels; the remaining entries are independently launched
kernels.  ``Staging'' and ``reduction'' entries implement parts of the cited
mathematical map; they do not define additional model equations.

\small
\begin{longtable}{>{\raggedright\arraybackslash}p{0.43\textwidth}
                  >{\raggedright\arraybackslash}p{0.51\textwidth}}
\toprule
Kernel & Mathematical role \\
\midrule
\endhead
\kernelref{pi_forward.py}{_load_event_log_probability} & Selects a family/species parameter layout and returns $\ell_{g,s}$. \\
\kernelref{pi_forward.py}{_compute_total_receiver_mass} & Computes $M_x$ and $R_x$ for Equation~\eqref{eq:transfer-complement}. \\
\kernelref{pi_forward.py}{_compute_transfer_complement} & Computes $\mathcal B[x](s)$ in Equation~\eqref{eq:transfer-complement}. \\
\kernelref{pi_forward.py}{_initialize_leaf_reconciliation_likelihood_kernel} & Leaf specialization of Equation~\eqref{eq:pi-wave}. \\
\kernelref{pi_forward.py}{_update_reconciliation_likelihood_kernel} & General primal recurrence, Equations~\eqref{eq:pi-wave} and \eqref{eq:pibar}. \\
\kernelref{pi_forward.py}{_reduce_single_gene_split_events_kernel} & Equation~\eqref{eq:dts} for one split per parent. \\
\kernelref{pi_forward.py}{_stage_multiple_gene_split_event_reduction_kernel} & Partial stabilized sums for Equation~\eqref{eq:dts} with multiple splits. \\
\kernelref{pi_forward.py}{_finalize_multiple_gene_split_event_reduction_kernel} & Final merge of the partial sums for Equation~\eqref{eq:dts}. \\
\kernelref{e_step.py}{_update_extinction_log_probabilities_kernel} & Extinction update in Equation~\eqref{eq:e-step}. \\
\kernelref{e_step.py}{_stage_extinction_and_transfer_complement_vjp_kernel} & Event adjoints from Equation~\eqref{eq:e-adjoint-event} and stages Equation~\eqref{eq:complement-vjp}. \\
\kernelref{e_step.py}{_finalize_extinction_transfer_complement_vjp_kernel} & Completes Equation~\eqref{eq:complement-vjp}. \\
\kernelref{e_step_tangent.py}{_update_extinction_log_probabilities_jvp_kernel} & Equations~\eqref{eq:lse-jvp} and \eqref{eq:ebar-jvp}. \\
\kernelref{e_step_so.py}{_stage_extinction_vjp_directional_derivative_kernel} & Directional derivative in Equation~\eqref{eq:e-so}, including Equation~\eqref{eq:quotient-tangent}. \\
\kernelref{e_step_so.py}{_finalize_extinction_vjp_directional_derivative_kernel} & Product-rule completion of the differentiated complementary-mass adjoint. \\
\kernelref{wave_tangent.py}{_update_reconciliation_likelihood_jvp_kernel} & JVP of Equations~\eqref{eq:pi-wave}--\eqref{eq:pibar}. \\
\kernelref{wave_tangent.py}{_apply_reconciliation_self_loop_jvp_iterations_kernel} & Repeated JVP of the within-row self-loop map. \\
\kernelref{wave_so.py}{_reconciliation_vjp_directional_derivative_kernel} & Wave second-order contraction in Equation~\eqref{eq:wave-so}. \\
\kernelref{dts_tangent.py}{_gene_split_reduction_jvp_kernel} & Gene-split JVP in Equation~\eqref{eq:dts-jvp}. \\
\kernelref{dts_so.py}{_gene_split_event_vjp_directional_derivative_kernel} & Directional derivative of the full-recurrence gene-split VJP, Equations~\eqref{eq:dts-full-probability} and \eqref{eq:dts-event-so}. \\
\kernelref{dts_so.py}{_transfer_subtree_vjp_directional_derivative_kernel} & Tree part of the transfer-complement second-order contraction, Equation~\eqref{eq:tree-complement-so}. \\
\kernelref{wave_backward_kernels.py}{_select_active_adjoint_rows_kernel} & Selects rows whose adjoint right-hand side exceeds the pruning threshold. \\
\kernelref{wave_backward_kernels.py}{_prepare_reconciliation_self_loop_vjp_kernel} & Forms event probabilities and coefficients of $J_{i,\mathrm{self}}^T$ for Equation~\eqref{eq:wave-adjoint-solve}. \\
\kernelref{wave_backward_kernels.py}{_apply_reconciliation_self_loop_transpose_kernel} & Applies $J_{i,\mathrm{self}}^T$ during the solve in Equation~\eqref{eq:wave-adjoint-solve}. \\
\kernelref{wave_backward_kernels.py}{_accumulate_transfer_receiver_log_probability_vjp_kernel} & Receiver-log-probability component of Equation~\eqref{eq:complement-vjp}. \\
\kernelref{wave_backward_kernels.py}{_accumulate_reconciliation_event_vjp_kernel} & Contracts the solved adjoint with the six named local-event probabilities. \\
\kernelref{wave_backward_kernels.py}{_accumulate_gene_split_event_vjp_kernel} & Full-recurrence gene-split VJP from Equation~\eqref{eq:dts-full-probability}: scatters event VJPs and stages transfer donor adjoints. \\
\kernelref{wave_backward_kernels.py}{_reduce_max_transfer_vjp_kernel} & Reduces cotangent partials of $m$ (\texttt{max\_transfer}) over split tiles. \\
\kernelref{wave_backward_kernels.py}{_select_active_transfer_donor_sides_kernel} & Selects nonzero staged donor-adjoint sides before Equation~\eqref{eq:tree-complement}. \\
\kernelref{wave_backward_kernels.py}{_accumulate_transfer_subtree_vjp_kernel} & Computes subtree sums and Equation~\eqref{eq:tree-complement}. \\
\bottomrule
\end{longtable}
\normalsize

\section{Synchronization and discrete layout contracts}

Tree reductions overwrite a state entry with its subtree sum.  A barrier is
required after computing the original total $A$: otherwise one warp can replace
$u(s)$ by $u(s)+u(c_1(s))+u(c_2(s))$ while another warp is still reducing $A$,
which double-counts the children.  A barrier is also required between tree
levels so a parent is not read before its children are complete.

Species-tree and split indices are discrete metadata, not numerical model
precision.  They are stored compactly as 32-bit integers.  Flattened address
products are widened to 64 bits where needed.  Parent-chain loops explicitly
retain a stable integer type because Triton requires loop-carried variables to
have the same type on every iteration.

\section{Numerical conventions}

All recurrences are evaluated in log2 space.  Max subtraction in log-sum-exp
and transfer complements is algebraically neutral.  Model residuals may be
float32 or float64; row offsets may use equal or wider accumulator precision.
An offset is narrowed only when a frame shift enters a residual-precision event
probability.  Split priors retained at accumulator precision are similarly
converted at the DTS residual boundary.

Atomic additions implement sums over gene splits, child scatters, families,
or receiver states.  Their ordering can change the last floating-point bits,
but not the mathematical sum described above.

\end{document}
